%%%%%%%%%%%%%%%%%%%%%%%%%%%%%%%%%%%%%%%%%%%%%%%%%%%%%%%%%%%%%%%%%%%%%%%%%%%%%
%  Langage graphique commun à toutes les figures.
%
%  Les treize figures du rapport partagent ces styles : une forme et une
%  couleur y ont partout le même sens. Le lecteur n'a donc à apprendre la
%  légende qu'une seule fois.
%
%     gris    donnée, fichier, artefact          rectangle
%     bleu    traitement déterministe (code)     rectangle
%     vert    modèle appris (embeddings, LLM)    rectangle
%     ambre   décision, garde-fou                losange
%     rouge   arrêt, abstention, échec           rectangle
%
%  Palette de données validée pour la vision des couleurs (deutan/protan/
%  tritan) : bleu, ambre, vert, prune. Le rouge est réservé au statut
%  « échec » et n'est jamais employé comme simple série.
%%%%%%%%%%%%%%%%%%%%%%%%%%%%%%%%%%%%%%%%%%%%%%%%%%%%%%%%%%%%%%%%%%%%%%%%%%%%%

% ============================================================== SCHÉMAS TikZ

\tikzset{
  % --- boîtes ---------------------------------------------------------------
  bloc/.style={
    draw, rounded corners=2.5pt, line width=0.7pt,
    minimum height=11mm, text width=23mm, align=center,
    inner sep=3pt, font=\sffamily\scriptsize},
  donnee/.style={bloc, draw=slate, fill=mist,          text=ink},
  traitement/.style={bloc, draw=inptblue, fill=inptblue!8,  text=navy},
  modele/.style={bloc, draw=leafgreen, fill=leafgreen!8, text=leafgreen!45!black},
  arret/.style={bloc, draw=anrtred, fill=anrtred!7,    text=anrtred!70!black},
  decision/.style={
    draw=amberdark, fill=amberdark!10, line width=0.7pt,
    diamond, aspect=2.1, align=center, inner sep=1pt,
    text width=17mm, font=\sffamily\scriptsize, text=amberdark!60!black},
  % variantes de gabarit
  large/.style={text width=30mm},
  etroit/.style={text width=18mm},
  haut/.style={minimum height=14mm},
  %
  % --- liens ----------------------------------------------------------------
  flux/.style={-{Stealth[length=2.2mm,width=1.6mm]}, line width=0.65pt, draw=slate},
  fluxfort/.style={flux, line width=1pt, draw=navy},
  usage/.style={-{Stealth[length=2.2mm,width=1.6mm]}, line width=0.55pt,
                draw=slate!70, dash pattern=on 2pt off 1.6pt},
  rejet/.style={-{Stealth[length=2.2mm,width=1.6mm]}, line width=0.65pt, draw=anrtred},
  %
  % --- textes ---------------------------------------------------------------
  etiq/.style={font=\sffamily\scriptsize, text=slate, inner sep=2pt},
  etiqfort/.style={font=\sffamily\scriptsize\bfseries, text=navy, inner sep=2pt},
  annot/.style={font=\sffamily\scriptsize\itshape, text=slate, align=left,
                inner sep=2pt},
  %
  % --- conteneurs -----------------------------------------------------------
  bande/.style={
    draw=rule, line width=0.7pt, rounded corners=3pt,
    dash pattern=on 3pt off 2pt, inner sep=6mm},
  bandepleine/.style={draw=rule, line width=0.7pt, rounded corners=3pt,
                      fill=mist!45, inner sep=6mm},
  titrebande/.style={
    font=\sffamily\scriptsize\bfseries, text=slate,
    anchor=west, fill=white, inner xsep=4pt, inner ysep=1pt},
  %
  % --- pastille de renvoi vers le texte -------------------------------------
  puce/.style={
    circle, draw=anrtred, fill=white, text=anrtred, line width=0.6pt,
    inner sep=0.6pt, minimum size=3.6mm, font=\sffamily\bfseries\tiny},
}

% Titre posé sur le bord supérieur gauche d'un conteneur.
% #1 = nœud conteneur, #2 = libellé
%
% Le décalage passe par « xshift » et non par la syntaxe calc « $...$ » :
% celle-ci ne survit pas à un \resizebox, qui réinterprète le signe dollar.
\newcommand{\bandetitre}[2]{%
  \node[titrebande, xshift=6mm] at (#1.north west) {#2};}

% Légende compacte des types de nœud, à poser sous un schéma.
\newcommand{\legendeschema}{%
  \begin{tikzpicture}[baseline]
    \node[donnee, text width=0pt, minimum width=5mm, minimum height=3.2mm] (a) {};
    \node[etiq, right=1.2mm of a] (la) {donnée};
    \node[traitement, text width=0pt, minimum width=5mm, minimum height=3.2mm,
          right=4mm of la] (b) {};
    \node[etiq, right=1.2mm of b] (lb) {traitement};
    \node[modele, text width=0pt, minimum width=5mm, minimum height=3.2mm,
          right=4mm of lb] (c) {};
    \node[etiq, right=1.2mm of c] (lc) {modèle appris};
    \node[decision, text width=0pt, minimum width=5mm, minimum height=2.2mm,
          right=4mm of lc] (d) {};
    \node[etiq, right=1.2mm of d] (ld) {décision};
    \node[arret, text width=0pt, minimum width=5mm, minimum height=3.2mm,
          right=4mm of ld] (e) {};
    \node[etiq, right=1.2mm of e] {arrêt};
  \end{tikzpicture}}

% ============================================================ FIGURES pgfplots

% Séparateur décimal français dans toutes les figures : le corps du texte
% écrit « 0,969 », les graphiques doivent en faire autant.
\pgfkeys{/pgf/number format/.cd, use comma, 1000 sep={\,}}

\pgfplotsset{
  % Socle commun : pas de cadre, une seule grille (celle de l'axe des
  % valeurs), graduations discrètes. Tout ce qui n'aide pas à lire un chiffre
  % est retiré.
  socle/.style={
    width=\linewidth, height=6.4cm,
    font=\sffamily\small,
    axis line style={draw=rule, line width=0.6pt},
    axis x line*=bottom, axis y line*=left,
    every axis label/.style={font=\sffamily\small, text=slate},
    tick label style={font=\sffamily\scriptsize, text=slate},
    label style={font=\sffamily\scriptsize, text=slate},
    title style={font=\sffamily\small\bfseries, text=navy, align=left},
    major tick length=2.5pt, minor tick length=0pt,
    tick align=outside,
    grid style={draw=rule!55, line width=0.4pt},
    legend style={
      draw=none, fill=none, font=\sffamily\scriptsize, text=slate,
      cells={anchor=west}, legend columns=-1,
      /tikz/every even column/.append style={column sep=8pt}},
    legend image post style={line width=2pt},
    clip=false,
  },
  % Barres horizontales : la catégorie se lit à plat, sans tête penchée.
  barresH/.style={
    socle,
    xbar, bar width=9pt,
    xmajorgrids, ymajorgrids=false,
    ytick=data, y dir=reverse,
    enlarge y limits=0.14,
    % L'axe vertical ne porte aucune quantité : on efface son trait et ses
    % graduations, mais surtout PAS « axis y line=none », qui emporterait
    % aussi les libellés de catégorie.
    axis x line*=bottom, axis y line*=left,
    y axis line style={draw=none},
    ytick style={draw=none},
    nodes near coords style={font=\sffamily\scriptsize, text=ink,
                             anchor=west, xshift=2pt},
  },
  % Barres verticales groupées.
  barresV/.style={
    socle,
    ybar, bar width=13pt,
    ymajorgrids, xmajorgrids=false,
    enlarge x limits=0.22,
    xtick=data,
    nodes near coords style={font=\sffamily\scriptsize, text=ink,
                             anchor=south, yshift=1pt},
  },
  % Nuage de points en bandes (scores d'abstention).
  bandesY/.style={
    socle,
    ymajorgrids=false, xmajorgrids=true,
    axis y line=none,
    enlarge y limits=0.3,
  },
}

% Cycle de couleurs des séries — ordre fixe, jamais recyclé.
\pgfplotscreateplotcyclelist{serie}{
  {draw=inptblue,  fill=inptblue},
  {draw=amberdark, fill=amberdark},
  {draw=leafgreen, fill=leafgreen},
  {draw=plum,      fill=plum},
}
\pgfplotsset{cycle list name=serie}

% Extrémité de barre arrondie côté valeur, ancrée sur la ligne de base.
\pgfplotsset{
  bout arrondi/.style={
    /pgfplots/every axis plot post/.append style={
      rounded corners=1.6pt}},
}

% Filet vertical de seuil, avec son étiquette.
% #1 = abscisse, #2 = libellé
\newcommand{\seuil}[2]{%
  \draw[anrtred, line width=0.8pt, dash pattern=on 3pt off 2pt]
    ({axis cs:#1,0}|-{rel axis cs:0,0}) -- ({axis cs:#1,0}|-{rel axis cs:0,1});
  \node[font=\sffamily\scriptsize\bfseries, text=anrtred, anchor=south, rotate=90]
    at ({axis cs:#1,0}|-{rel axis cs:0,0.02}) {\;#2};}

% ================================================================ DIAGRAMME
% Gantt : mêmes couleurs, mêmes fontes que le reste.
\ganttset{
  bar height=0.55,
  bar/.append style={draw=inptblue!70, fill=inptblue!25, rounded corners=1pt},
  bar label font=\sffamily\scriptsize\color{ink},
  group/.append style={draw=navy, fill=navy},
  group label font=\sffamily\scriptsize\bfseries\color{navy},
  milestone/.append style={draw=amberdark, fill=amberdark},
  milestone label font=\sffamily\scriptsize\color{amberdark!60!black},
  title label font=\sffamily\scriptsize\color{slate},
  title/.append style={draw=none, fill=none},
  link/.style={-{Stealth[length=1.6mm]}, draw=slate, line width=0.5pt},
  vgrid={*1{draw=rule!60, line width=0.4pt}},
}
